\documentclass[11pt]{article}
\usepackage[margin=1in]{geometry}
\usepackage{amsmath}
\usepackage{booktabs}

\title{Naive Bayes Revisited: Competitive Digit Recognition Without Training Loops}
\author{Paper2Proof canary set\thanks{A small paper written for the Paper2Proof evaluation set (eval set A).}}
\date{2026}

\begin{document}
\maketitle

\begin{abstract}
Gaussian naive Bayes fits in a single pass over the data and has no hyperparameters to tune. We
show that on the scikit-learn digits dataset it reaches a test accuracy of 0.920, close to
iteratively trained models, at a fraction of their cost.
\end{abstract}

\section{Introduction}
Closed-form classifiers remain attractive when compute is scarce. We evaluate Gaussian naive
Bayes on handwritten digits under a fixed, fully specified protocol.

\section{Method}
Gaussian naive Bayes with one diagonal Gaussian per class, fitted by maximum likelihood on the
raw pixel values, with the default variance smoothing of $10^{-9}$ times the largest feature
variance. No feature scaling, no feature selection.

\section{Experimental setup}
\paragraph{Data.} The digits dataset shipped with scikit-learn (\texttt{load\_digits}): 1797
images of $8 \times 8$ pixels, 10 classes.
\paragraph{Splits.} For each seed $s \in \{0, 1, 2\}$, a stratified split with 30\,\% test
samples (\texttt{train\_test\_split(test\_size=0.3, random\_state=s, stratify=y)}).
\paragraph{Metric.} Test accuracy, mean and sample standard deviation over the three seeds.
\paragraph{Software.} Python 3.12, scikit-learn 1.9.1, CPU only.

\section{Results}
As shown in Table~\ref{tab:digits-nb}, Gaussian naive Bayes reaches a test accuracy of
0.920 $\pm$ 0.010 on digits.

\begin{table}[t]
\centering
\caption{Test accuracy on digits (mean $\pm$ standard deviation over seeds 0, 1, 2).}
\label{tab:digits-nb}
\begin{tabular}{lc}
\toprule
Model & Accuracy \\
\midrule
Gaussian naive Bayes & $0.920 \pm 0.010$ \\
\bottomrule
\end{tabular}
\end{table}

\section{Conclusion}
Naive Bayes is a competitive, training-free classifier for small digit recognition tasks.

\end{document}
